% Bab V  Kesimpulan dan Saran
\chapter{KESIMPULAN DAN SARAN}

\section{Kesimpulan}
Kesimpulan merupakan uraian singkat hasil penelitian dan pembahasan (bukan kalimat statistik) yang mengacu pada tujuan penelitian. Tuliskan sesuai urutan rumusan masalah. Kesimpulan dapat diringkas dalam tabel yang menghubungkannya dengan rumusan masalah, seperti Tabel 5.1 (isi tabel adalah contoh dan tidak memuat hal baru di luar Bab IV).

\begin{singlespace}
\begin{longtable}{|c|p{0.36\linewidth}|p{0.46\linewidth}|}
\caption{Contoh ringkasan kesimpulan terhadap rumusan masalah}\\
\hline
\textbf{No.} & \textbf{Rumusan masalah} & \textbf{Kesimpulan} \\
\hline
1 & Bagaimana merancang sistem pemantauan suhu berbasis ESP32? & Sistem berhasil dirancang dan berfungsi pada seluruh pengujian (Gambar 4.1). \\
\hline
2 & Seberapa besar akurasi sensor suhu yang digunakan? & Galat rata-rata sensor 1,25\%, memenuhi indikator $\leq$ 2\% (Tabel 4.2). \\
\hline
\end{longtable}
\end{singlespace}

\section{Saran}
Saran berupa anjuran yang mengacu pada kesimpulan dan dapat dimanfaatkan oleh pihak yang berkepentingan, termasuk arah pengembangan penelitian berikutnya.

\textit{Contoh kalimat bersitasi:} Penelitian selanjutnya dapat merujuk pada kebijakan terbaru (Kementerian Komunikasi dan Digital, 2025).
